\documentclass{article}
\usepackage{iclr2026_conference,times}

% ─── Packages ──────────────────────────────────────────────────────────────
\usepackage{amsmath, amssymb, amsthm}
\usepackage{bm}
\usepackage{graphicx}
\usepackage{booktabs}
\usepackage{multirow}
\usepackage{array}
\usepackage{tabularx}
\usepackage{algorithm}
\usepackage{algpseudocode}
\usepackage{enumitem}
\usepackage{tikz}
\usetikzlibrary{arrows.meta, positioning, shapes.geometric, fit, backgrounds}
\usepackage{hyperref}
\usepackage{url}
\usepackage{xcolor}
\usepackage{xspace}

% ─── Theorem environments ───────────────────────────────────────────────────
\newtheorem{definition}{Definition}[section]
\newtheorem{theorem}{Theorem}[section]
\newtheorem{proposition}{Proposition}[section]
\newtheorem{remark}{Remark}[section]
\newtheorem{lemma}{Lemma}[section]

% ─── Custom commands ────────────────────────────────────────────────────────
\newcommand{\memset}{\mathcal{M}}
\newcommand{\note}{m}
\newcommand{\ours}{\textsc{ASEM}\xspace}
\newcommand{\ie}{\textit{i.e.},\xspace}
\newcommand{\eg}{\textit{e.g.},\xspace}
\newcommand{\cf}{\textit{cf.}\xspace}
\newcommand{\etal}{\textit{et al.}\xspace}

% ─── Title ──────────────────────────────────────────────────────────────────
\title{\ours: Agentic Self-Evolving Memory for Large Language Model Agents\\
\large Structured Organisation, Reinforcement-Learned Operations, and Value-Aware Retrieval}

\author{Anonymous authors}

% ════════════════════════════════════════════════════════════════════════════
\begin{document}

\maketitle

% ════════════════════════════════════════════════════════════════════════════
\begin{abstract}
Large language model (LLM) agents exhibit remarkable competence in isolated
tasks yet remain fundamentally \emph{stateless}: information encountered in
one session is irrecoverably lost before the next, imposing a structural
ceiling on long-horizon adaptation.
We identify three root causes of this ceiling.
\textbf{(i) Semantic impoverishment:} raw interaction logs stored without
enrichment produce flat, noise-ridden memory banks whose retrieval is
dominated by surface-form similarity rather than conceptual relevance.
\textbf{(ii) Heuristic write policies:} append-only or recency-based eviction
strategies fragment complementary information and accumulate unresolved
contradictions, causing retrieval quality to degrade monotonically with
corpus growth.
\textbf{(iii) Utility blindness:} policies that equate semantic proximity with
functional value return contextually close but strategically inferior
candidates, systematically hindering multi-hop reasoning and long-horizon task
completion.

We present \ours (\textbf{A}gentic \textbf{S}elf-\textbf{E}volving
\textbf{M}emory), a unified five-stage framework that addresses all three
deficiencies simultaneously.
\ours introduces: a \emph{multi-attribute atomic note schema} that encodes
each experience along keyword, categorical, contextual, embedding, and utility
axes; an \emph{RL-trained Memory Manager} that selects among four principled
write operations using outcome-based reward rather than surface heuristics;
a \emph{dynamic linking and bidirectional memory evolution} mechanism that
weaves newly integrated notes into a traversable knowledge network while
retroactively revising the representations of related prior entries; a
\emph{two-phase hybrid retrieval pipeline} that decouples semantic relevance
from learned functional utility, followed by a distillation step that filters
retrieved candidates to the minimal noise-free subset; and a
\emph{non-parametric runtime utility estimation} rule whose convergence in
expectation to the true per-memory expected return, with bounded asymptotic
variance, is formally guaranteed.
The resulting system maintains a living knowledge network that
self-organises, self-repairs, and continuously refines utility estimates of
stored experiences---without ever modifying backbone LLM parameters.
Empirical evaluation on long-horizon conversational benchmarks demonstrates
substantial gains over strong single-paradigm baselines across all
capability dimensions.
\end{abstract}


% ════════════════════════════════════════════════════════════════════════════
\section{Introduction}
\label{sec:introduction}
% ════════════════════════════════════════════════════════════════════════════

\subsection{The Statelesness Problem}
\label{subsec:statelesness}

The deployment of LLM-powered agents in real-world, multi-session environments
exposes a fundamental architectural mismatch.
These agents operate under a hard context-window bound: every interaction
begins from a blank slate, with no access to experiential knowledge accumulated
in prior sessions.
For tasks whose successful execution depends on information that spans multiple
interactions---tracking evolving user preferences, resolving references to
entities mentioned weeks earlier, or building a coherent model of a
long-running project---this statelesness is not a minor inconvenience but a
structural barrier~\cite{wu2024longmemeval, maharana2024locomo}.

Human cognition resolves an analogous problem through the interplay of
\emph{episodic} and \emph{semantic} memory~\cite{tulving1972episodic}.
Episodic memory records specific past experiences; semantic memory distils
those episodes into generalised knowledge structures; and a constructive
retrieval process synthesises relevant fragments from both sources into
actionable representations tailored to the current demand.
Crucially, this process is \emph{selective}: not every stored trace is
retrieved, and retrieved traces are not replayed verbatim but
\emph{reconstructed} in accordance with the current context~\cite{schacter2007constructive}.

Replicating this capacity in LLM agents demands solving four interlocking
sub-problems:
\textbf{(P1)} constructing semantically rich memory representations that
capture the \emph{implicit} structure of an interaction rather than its
surface tokens;
\textbf{(P2)} executing principled write operations that consolidate
complementary facts and resolve contradictions rather than blindly accumulating
entries;
\textbf{(P3)} maintaining relational structures that connect semantically and
causally related notes into a traversable knowledge network for multi-hop
reasoning; and
\textbf{(P4)} retrieving memories that are functionally \emph{useful} for a
query---a property strictly stronger than semantic \emph{relevance} that
cosine similarity cannot capture.

\subsection{The Integration Gap}
\label{subsec:gap}

The literature on memory-augmented LLM agents partitions into three paradigms,
each addressing a non-overlapping subset of $\{\text{P1}, \text{P2}, \text{P3},
\text{P4}\}$.

\paragraph{Static storage and retrieval.}
Systems that augment agents with plain key-value memory banks retrieved by
embedding similarity are effective for isolated factual recall, partially
addressing P1.
They provide no mechanism for memory organisation (P2), relational linking
(P3), or utility-aware selection (P4).
Retrieval quality degrades monotonically with corpus growth as contradictory
and redundant entries accumulate~\cite{packer2023memgpt, zhong2024memorybank}.

\paragraph{Structured and graph-based memory.}
Systems that incorporate predefined schemas and typed relational structures
partially address P3 for anticipated relationship types.
Their reliance on fixed schemas fundamentally limits adaptability: novel
relationship types encountered at runtime cannot be accommodated without
schema modification, causing memory organisation to converge to a static
predefined structure rather than emerging organically from
content~\cite{chhikara2025mem0}.

\paragraph{Reinforcement-learning-enhanced memory.}
Systems that train dedicated agents to perform structured write operations
address P2, and when combined with utility estimation, partially address P4.
However, existing implementations either operate on semantically impoverished
flat entries, failing P1 and P3, or provide principled utility estimation
without the rich relational structure that enables multi-hop reasoning,
failing P3~\cite{memoryr12025, memrl2026}.

No existing system simultaneously satisfies all four requirements.
We term this the \textbf{integration gap} and present \ours as its resolution.

\subsection{Contributions}
\label{subsec:contributions}

\begin{enumerate}[leftmargin=1.8em, label=\textbf{(\arabic*)}]

  \item \textbf{Unified multi-attribute note representation (P1).}
  We introduce a structured note schema that enriches each memory entry with
  LLM-generated keywords, categorical tags, a contextual description, a dense
  embedding, an intent embedding, bidirectional link sets, and a scalar
  utility (Q-value). This multi-faceted structure supports simultaneous
  retrieval along keyword, topic, contextual, and utility axes.

  \item \textbf{RL-driven principled memory management (P2).}
  We train a Memory Manager using Group Relative Policy Optimisation
  (GRPO)~\cite{shao2024deepseekmath} with outcome-based reward derived from a
  frozen Answer Agent, teaching the manager to consolidate complementary facts
  via \textsc{update} operations rather than fragmenting the memory
  space---a failure mode that systematically afflicts heuristic managers.

  \item \textbf{Dynamic linking and bidirectional memory evolution (P3).}
  We design a mechanism by which newly integrated memories trigger simultaneous
  link generation into the existing knowledge network and retroactive
  contextual revision of related existing notes, enabling higher-order
  conceptual patterns to emerge over time.

  \item \textbf{Value-aware two-phase retrieval with distillation (P4).}
  We propose a pipeline that decouples semantic relevance from functional
  utility. Phase~A constructs a candidate pool via similarity filtering;
  Phase~B re-ranks by a composite score interpolating z-score-normalised
  similarity and learned Q-value. A downstream Answer Agent then distils
  the retrieved set to the minimal noise-free subset.

  \item \textbf{Formal convergence guarantees.}
  We prove that the EMA-based utility update rule converges in expectation to
  the true per-memory expected return with bounded asymptotic variance, and
  that the composite retrieval policy is optimal under a KL-regularised
  variational lower bound on expected reward.

\end{enumerate}


% ════════════════════════════════════════════════════════════════════════════
\section{Related Work}
\label{sec:related}
% ════════════════════════════════════════════════════════════════════════════

Table~\ref{tab:related_comparison} provides a structured comparison across
five capability dimensions corresponding directly to sub-problems P1--P4 and
runtime evolution capability.

\begin{table}[t]
\centering
\caption{Structured comparison of closely related memory systems.
\textbf{Rich}: rich note representation (P1);
\textbf{Write}: principled learned write operations (P2);
\textbf{Link}: dynamic relational linking (P3);
\textbf{Util}: utility-aware retrieval (P4);
\textbf{Evolve}: non-parametric runtime self-evolution.
\checkmark~fully supported; $\circ$~partially; $\times$~absent.}
\label{tab:related_comparison}
\vspace{4pt}
\renewcommand{\arraystretch}{1.25}
\begin{tabularx}{\linewidth}{l l c c c c c}
\toprule
\textbf{System} & \textbf{Paradigm} & \textbf{Rich} & \textbf{Write} &
\textbf{Link} & \textbf{Util} & \textbf{Evolve} \\
\midrule
MemGPT~\cite{packer2023memgpt}        & OS-inspired hierarchy    & $\times$ & $\circ$ & $\times$ & $\times$ & $\times$ \\
MemoryBank~\cite{zhong2024memorybank} & Forgetting-curve store   & $\times$ & $\circ$ & $\times$ & $\times$ & $\times$ \\
ReadAgent~\cite{lee2024readagent}     & Gist distillation        & $\circ$  & $\times$& $\times$ & $\times$ & $\times$ \\
Mem0~\cite{chhikara2025mem0}          & Graph + LLM ops          & $\circ$  & $\circ$ & $\circ$  & $\times$ & $\times$ \\
MemoryOS~\cite{kang2025memoryos}      & OS-abstraction tiers     & $\circ$  & $\circ$ & $\times$ & $\times$ & $\times$ \\
Atomic-Linking~\cite{amem2025}        & Zettelkasten-inspired    & \checkmark & $\times$ & \checkmark & $\times$ & $\circ$ \\
RL-Manager~\cite{memoryr12025}        & RL write operations      & $\times$ & \checkmark & $\times$ & $\times$ & $\times$ \\
Value-Retrieval~\cite{memrl2026}      & Q-value retrieval        & $\circ$  & $\times$& $\times$ & \checkmark & \checkmark \\
\midrule
\textbf{\ours (ours)} & Unified & \checkmark & \checkmark & \checkmark & \checkmark & \checkmark \\
\bottomrule
\end{tabularx}
\end{table}

\subsection{Memory Systems for LLM Agents}

The earliest deployed approach maintains complete interaction histories passed
as prefix context~\cite{maharana2024locomo} or compressed by a reading
agent~\cite{lee2024readagent} to fit within the context window.
While effective for recent events, these systems scale poorly with session
length and provide no mechanism for selectively surfacing the most relevant
historical information.

\paragraph{Flat key-value stores.}
MemGPT~\cite{packer2023memgpt} draws inspiration from OS memory hierarchies to
implement a two-tier main-context / external-context architecture with
rudimentary CRUD operations.
Operation selection is governed by pre-trained LLM priors rather than
task-outcome-optimised policies, leaving the system prone to consolidation
failures: complementary facts are fragmented across entries and contradictions
coexist indefinitely.
MemoryBank~\cite{zhong2024memorybank} augments flat storage with an
Ebbinghaus-inspired forgetting curve that attenuates rarely retrieved entries
while reinforcing frequently accessed ones---an elegant macroscopic
consolidation mechanism that nonetheless operates on raw text embeddings
without semantic enrichment or relational linking.

\paragraph{Graph-enhanced systems.}
Mem0~\cite{chhikara2025mem0} extends the CRUD operation set with
graph-database-backed storage and relationship modelling.
Operation selection remains governed by vanilla LLM prompting without any
learning signal, preserving consolidation failure modes.
MemoryOS~\cite{kang2025memoryos} frames memory management as an OS abstraction
with hot and cold storage tiers but relies on heuristic promotion and eviction
policies without downstream task signal.

\paragraph{Atomic note construction and dynamic linking.}
A Zettelkasten-inspired approach~\cite{amem2025} operationalises atomic note
construction and dynamic inter-note linking in a computational memory system.
Comprehensive notes comprising raw content, keywords, categorical tags,
contextual descriptions, dense embeddings, and semantic link sets are
constructed upon arrival.
When a new note is integrated, its $k$ nearest neighbours in embedding space
are identified, and an LLM establishes meaningful connections (\emph{link
generation}) and retroactively updates neighbours' representations
(\emph{memory evolution}).
Ablation studies confirm that both operations are critical: removing either
significantly degrades multi-hop reasoning.
The limitation is the complete absence of principled write-operation selection
and any measure of functional utility.

\paragraph{RL-trained write operations.}
A reinforcement learning framework for memory management~\cite{memoryr12025}
trains a Memory Manager selecting among
$\{\textsc{add}, \textsc{update}, \textsc{delete}, \textsc{noop}\}$ and an
Answer Agent applying Memory Distillation, jointly trained with GRPO using
exact-match correctness as reward.
The contrast with vanilla LLM managers is compelling: faced with sequential
statements ``I adopted a dog named Buddy'' and ``I adopted another dog named
Scout'', a vanilla manager issues \textsc{delete}+\textsc{add} (fragmentation),
while the RL-trained manager issues a single \textsc{update} (consolidation).
The framework operates on flat, semantically impoverished entries and retrieves
by similarity alone.

\paragraph{Value-aware non-parametric retrieval.}
A non-parametric retrieval framework~\cite{memrl2026} organises memory as
Intent-Experience-Utility triplets and performs two-phase retrieval: Phase~A
filters by semantic similarity to a threshold; Phase~B re-ranks using a
composite score interpolating normalised similarity and a learned Q-value.
Utility estimates are updated via temporal-difference rules after each
interaction.
The principal limitation is the absence of rich relational structure between
memory entries, precluding multi-hop traversal.

\subsection{Retrieval-Augmented Generation}

Standard retrieval-augmented generation~\cite{lewis2020rag} indexes corpora
into dense vector chunks and retrieves the top-$k$ most similar at query time.
Advanced variants introduce query rewriting, hypothetical document embeddings,
and post-retrieval re-ranking~\cite{gao2023rag_survey}.
Agentic extensions treat retrieval as dynamic and iterative, interleaving
retrieval with chain-of-thought generation~\cite{asai2023selfrag,
trivedi2022interleaving}.
These approaches exhibit agency at the retrieval phase but maintain static
knowledge bases that do not evolve in response to interaction experience.
\ours differs fundamentally: its memory bank is dynamically written, linked,
and evolved by agent interactions, and retrieval incorporates learned utility
estimates that reflect demonstrated functional value.

\subsection{Continual and Lifelong Learning}

Classical continual learning approaches~\cite{kirkpatrick2017ewc,
parisi2019continual} mitigate catastrophic forgetting by constraining
parameter updates or preserving samples from past distributions.
For LLM-scale models, parametric adaptation is computationally prohibitive and
risks destabilising pre-trained alignments.
This motivates non-parametric continual learning, where the backbone is frozen
and adaptation occurs through modifications to an external memory
structure~\cite{memrl2026}.
\ours squarely adopts this paradigm: the backbone LLM is never modified.
Runtime adaptation occurs entirely through memory writes, linkages, evolutions,
and scalar Q-value updates.


% ════════════════════════════════════════════════════════════════════════════
\section{Notation and Formal Problem Setting}
\label{sec:preliminaries}
% ════════════════════════════════════════════════════════════════════════════

\subsection{Memory Bank and Note Structure}

Let $\memset_t = \{\note_i\}_{i=1}^{N_t}$ denote the memory bank at discrete
time step $t$.
Each memory note is a structured tuple
\begin{equation}
  \note_i \;=\;
  \bigl(c_i,\; t_i,\; \mathcal{K}_i,\; \mathcal{G}_i,\; X_i,\;
        \mathbf{e}_i,\; \mathcal{L}_i,\; \mathbf{z}_i,\; q_i\bigr),
  \label{eq:note}
\end{equation}
where $c_i$ is the raw interaction content; $t_i$ is a timestamp;
$\mathcal{K}_i$, $\mathcal{G}_i$, and $X_i$ are LLM-generated keywords,
categorical tags, and a contextual description, respectively;
$\mathbf{e}_i \in \mathbb{R}^d$ is a dense embedding computed by encoding all
textual fields jointly; $\mathcal{L}_i \subseteq \memset_t$ is the set of
semantically linked notes; $\mathbf{z}_i \in \mathbb{R}^d$ is an intent
embedding representing the query context under which the note was created; and
$q_i \equiv q(\mathbf{z}_i, \note_i) \in \mathbb{R}$ is a learned utility
(Q-value) reflecting the expected downstream task reward associated with
retrieving $\note_i$ for queries with intent similar to $\mathbf{z}_i$.

\subsection{Memory-Based Markov Decision Process}

We model agent operation as a \emph{Memory-Based Markov Decision Process}
(M-MDP) $\langle \mathcal{S}, \mathcal{A}, \mathcal{P}, \mathcal{R}, \gamma,
\memset \rangle$, where the agent policy decomposes as
\begin{equation}
  \pi(y_t \mid s_t, \memset_t)
  \;=\;
  \sum_{\note \in \memset_t}
  \underbrace{\mu(\note \mid s_t, \memset_t)}_{\text{retrieval policy}}
  \cdot
  \underbrace{p_{\mathrm{LLM}}(y_t \mid s_t, \note)}_{\text{inference policy}},
  \label{eq:policy_decomp}
\end{equation}
where $\mu$ is the retrieval policy to be optimised and $p_{\mathrm{LLM}}$
is held frozen.
The central objective of \ours is to learn an optimal retrieval policy
\begin{equation}
  \mu^*
  \;=\;
  \operatorname*{arg\,max}_{\mu}\;
  \mathbb{E}\!\left[\sum_{t=0}^{\infty} \gamma^t\, R(s_t, y_t) \;\Big|\; \mu\right]
  \label{eq:objective}
\end{equation}
while maintaining structural coherence across $\memset_t$ and without
modifying the backbone parameters of $p_{\mathrm{LLM}}$.

\subsection{The Four Sub-Problems Formalised}

\begin{description}[leftmargin=1.5em, style=nextline]
  \item[\textbf{P1 --- Semantic enrichment.}]
  Given raw content $c$, construct $\note$ such that retrieval by any of
  $\mathbf{e}$, $\mathcal{K}$, $\mathcal{G}$, $X$ yields conceptually
  relevant results superior to retrieval by $c$ alone.
  \item[\textbf{P2 --- Principled write selection.}]
  Given new information $x$ and current bank $\memset$, select operation
  $o \in \{\textsc{add}, \textsc{update}, \textsc{delete}, \textsc{noop}\}$
  and target $\note'$ to maximise the subsequent task-conditional reward.
  \item[\textbf{P3 --- Relational structure.}]
  Maintain a non-trivial link graph $\mathcal{L} = \{\mathcal{L}_i\}$ such
  that multi-hop traversal over $\mathcal{L}$ recovers information not
  recoverable by single-shot similarity search.
  \item[\textbf{P4 --- Utility-aware retrieval.}]
  Retrieve the subset $\memset_{\mathrm{ret}} \subseteq \memset$ that
  maximises expected task reward, not merely semantic proximity to the query
  embedding.
\end{description}


% ════════════════════════════════════════════════════════════════════════════
\section{The \ours Framework}
\label{sec:methodology}
% ════════════════════════════════════════════════════════════════════════════

\ours comprises five tightly coupled stages.
Algorithm~\ref{alg:asem_overview} presents the complete system-level pipeline;
the individual stages are detailed below.

% ────────────────────────────────────────────────────────────────────────────
\subsection{Stage 1: Multi-Attribute Atomic Note Construction}
\label{sec:stage1}

\paragraph{Motivation.}
A memory entry stored as raw interaction tokens is simultaneously redundant
(retaining every filler word) and impoverished (capturing no latent semantic
structure).
Retrieval over such entries is dominated by lexical overlap, and merging
related entries into coherent representations is impossible without semantic
metadata.
We address P1 by constructing each note as a multi-attribute object whose
components provide orthogonal handles for downstream retrieval and linking.

\paragraph{Procedure.}
When a new interaction $c_{\mathrm{new}}$ arrives at time $t_{\mathrm{new}}$,
the backbone LLM is invoked with a structured prompt $P_1$ to generate
semantic attributes:
\begin{equation}
  \mathcal{K}_{\mathrm{new}},\; \mathcal{G}_{\mathrm{new}},\; X_{\mathrm{new}}
  \;\leftarrow\;
  \mathrm{LLM}\!\left(c_{\mathrm{new}} \;\|\; t_{\mathrm{new}} \;\|\; P_1\right).
  \label{eq:note_construction}
\end{equation}
A dense embedding is computed by encoding all textual fields jointly:
\begin{equation}
  \mathbf{e}_{\mathrm{new}}
  \;=\;
  f_{\mathrm{enc}}\!\Bigl[\mathrm{concat}\bigl(
    c_{\mathrm{new}},\; \mathcal{K}_{\mathrm{new}},\;
    \mathcal{G}_{\mathrm{new}},\; X_{\mathrm{new}}\bigr)\Bigr].
  \label{eq:embedding}
\end{equation}
The note is further initialised with an intent embedding
$\mathbf{z}_{\mathrm{new}} = f_{\mathrm{enc}}(c_{\mathrm{new}})$ and a prior
utility $q_{\mathrm{new}} = q_0$ that will be refined through the utility
update mechanism in Stage~5.
The complete note is then
\begin{equation}
  \note_{\mathrm{new}}
  \;=\;
  \bigl(c_{\mathrm{new}},\; t_{\mathrm{new}},\;
        \mathcal{K}_{\mathrm{new}},\; \mathcal{G}_{\mathrm{new}},\;
        X_{\mathrm{new}},\; \mathbf{e}_{\mathrm{new}},\;
        \emptyset,\; \mathbf{z}_{\mathrm{new}},\; q_0\bigr).
  \label{eq:new_note}
\end{equation}

\paragraph{Design rationale.}
The tripartite semantic enrichment $(\mathcal{K}, \mathcal{G}, X)$ creates
orthogonal representations of the same knowledge unit: keywords support
exact-match retrieval; categorical tags support topic-level filtering;
contextual descriptions support semantic similarity at a conceptual rather
than lexical level.
The joint encoding in Eq.~\eqref{eq:embedding} ensures that similarity search
captures the full semantic footprint of the note, including implicit structure
not present in the raw tokens.

% ────────────────────────────────────────────────────────────────────────────
\subsection{Stage 2: RL-Driven Memory Write Operations}
\label{sec:stage2}

\paragraph{Motivation.}
Naïve append-only policies accumulate redundant and contradictory entries
monotonically, degrading retrieval precision over time.
Recency-based eviction destroys potentially critical historical knowledge.
What is required is a policy that can recognise when a new piece of
information \emph{complements} an existing entry (warranting \textsc{update}),
\emph{contradicts} it (warranting \textsc{delete}), is \emph{genuinely new}
(warranting \textsc{add}), or is \emph{already fully represented} (warranting
\textsc{noop}).
We address P2 by training a dedicated Memory Manager to make this
discrimination using downstream task performance as the sole training signal.

\paragraph{Operation taxonomy.}
Let $\memset_{\mathrm{old}}$ be the current memory bank and $x$ the
information extracted from the new interaction.
The Memory Manager $\pi_\theta$ selects an operation $o$ and, where
applicable, a target entry $\note'$:
\begin{equation}
  (o,\; \note')
  \;\sim\;
  \pi_\theta\!\left(\cdot \;\middle|\; x,\; \memset_{\mathrm{old}}\right),
  \qquad
  o \in \{\textsc{add},\; \textsc{update},\; \textsc{delete},\; \textsc{noop}\}.
  \label{eq:memory_manager}
\end{equation}

\paragraph{Training objective.}
The Memory Manager is fine-tuned using Group Relative Policy Optimisation
(GRPO)~\cite{shao2024deepseekmath}, which samples $G$ candidate operations per
state and normalises advantages within the group to eliminate an explicit value
baseline:
\begin{equation}
  \mathcal{J}(\theta)
  \;=\;
  \mathbb{E}\!\left[
    \frac{1}{G}\sum_{i=1}^{G} \rho_\theta^{(i)}\, A_i
    \;-\;
    \beta\, D_{\mathrm{KL}}\!\left[\pi_\theta \;\|\; \pi_{\mathrm{ref}}\right]
  \right],
  \label{eq:grpo}
\end{equation}
where
$\rho_\theta^{(i)} = \pi_\theta(o^{(i)}, \note'^{(i)} \mid x,
\memset_{\mathrm{old}}) / \pi_{\mathrm{old}}(o^{(i)}, \note'^{(i)} \mid x,
\memset_{\mathrm{old}})$
is the per-action importance ratio and
$A_i = (r_i - \bar{r}) / \mathrm{std}(\mathbf{r})$
is the group-relative advantage.
The reward $r_i$ is the exact-match correctness of a \emph{frozen} Answer
Agent evaluated on the memory bank resulting from operation $o^{(i)}$:
\begin{equation}
  R_{\mathrm{answer}}
  \;=\;
  \mathrm{EM}(y_{\mathrm{pred}},\; y_{\mathrm{gold}}).
  \label{eq:reward}
\end{equation}
This outcome-driven signal requires no manual operation labels and teaches
the manager to consolidate complementary facts via \textsc{update} rather
than naïvely appending them.

% ────────────────────────────────────────────────────────────────────────────
\subsection{Stage 3: Dynamic Linking and Bidirectional Memory Evolution}
\label{sec:stage3}

\paragraph{Motivation.}
Even a well-organised flat store cannot support multi-hop reasoning unless
related notes are explicitly cross-referenced.
We address P3 by weaving newly integrated notes into the existing knowledge
network and retroactively revising the representations of related existing
notes to reflect the new information.

\paragraph{Link generation.}
After the write operation in Stage~2, the system identifies the $k$ nearest
neighbours of $\note_{\mathrm{new}}$ in $\memset$ by cosine embedding
similarity:
\begin{equation}
  s_{n,j} \;=\;
  \frac{\mathbf{e}_{\mathrm{new}} \cdot \mathbf{e}_j}
       {\|\mathbf{e}_{\mathrm{new}}\|\,\|\mathbf{e}_j\|},
  \qquad
  \memset_{\mathrm{new}}^{\mathrm{near}}
  \;=\;
  \bigl\{\note_j \mid \mathrm{rank}(s_{n,j}) \leq k,\;
  \note_j \in \memset\bigr\}.
  \label{eq:nearest_neighbours}
\end{equation}
The backbone LLM then analyses semantic and causal connections among
$\note_{\mathrm{new}}$ and $\memset_{\mathrm{new}}^{\mathrm{near}}$ to
populate the link set:
\begin{equation}
  \mathcal{L}_{\mathrm{new}}
  \;\leftarrow\;
  \mathrm{LLM}\!\left(
    \note_{\mathrm{new}} \;\|\;
    \memset_{\mathrm{new}}^{\mathrm{near}} \;\|\; P_2
  \right).
  \label{eq:link_generation}
\end{equation}
Links are generated without a predefined schema: the LLM produces free-form
relationship descriptions that capture semantic, causal, temporal, or any
other cognitively salient connection that the content warrants, permitting the
knowledge network to develop richer relational patterns as the agent
accumulates experience.

\paragraph{Bidirectional memory evolution.}
New information may render the contextual descriptions or categorical tags of
existing notes stale or incomplete.
For each $\note_j \in \memset_{\mathrm{new}}^{\mathrm{near}}$, the system
conditionally updates its attributes to incorporate the new information:
\begin{equation}
  \note_j^*
  \;\leftarrow\;
  \mathrm{LLM}\!\left(
    \note_{\mathrm{new}} \;\|\;
    \memset_{\mathrm{new}}^{\mathrm{near}} \setminus \{\note_j\} \;\|\;
    \note_j \;\|\; P_3
  \right),
  \label{eq:memory_evolution}
\end{equation}
and $\note_j^*$ replaces $\note_j$ in $\memset$.
This bidirectional update mechanism mirrors human memory reconsolidation: the
arrival of a new experience does not merely extend the store but
\emph{revises} the contextual framing of related prior knowledge, enabling the
emergence of higher-order conceptual patterns that no individual episode would
have produced in isolation.

% ────────────────────────────────────────────────────────────────────────────
\subsection{Stage 4: Two-Phase Hybrid Retrieval with Memory Distillation}
\label{sec:stage4}

\paragraph{Motivation.}
At inference time, the system must extract a compact, noise-free context from
$\memset$ for the frozen backbone.
Purely semantic retrieval conflates relevance with utility: a note
semantically close to the query may encode a strategy that failed in
analogous past tasks.
Conversely, a note whose intent embedding is distant but whose Q-value is
high may encode the structural pattern critical for the current multi-step
reasoning chain.
We address P4 by separating retrieval into two functionally distinct phases
followed by a learned distillation step.

\paragraph{Phase A --- Similarity-based candidate recall.}
Given a query $\mathbf{q}$, a candidate pool of contextually consistent notes
is assembled:
\begin{equation}
  \mathcal{C}(\mathbf{q})
  \;=\;
  \mathrm{TopK}_{k_1}\!\left(
    \bigl\{i \mid \mathrm{sim}(\mathbf{e}_\mathbf{q},\, \mathbf{e}_i) > \delta\bigr\},\;
    \text{by sim}
  \right),
  \label{eq:phase_a}
\end{equation}
where $\mathbf{e}_\mathbf{q} = f_{\mathrm{enc}}(\mathbf{q})$ and $\delta$ is
a sparsity threshold.
If $\mathcal{C}(\mathbf{q}) = \emptyset$, the system falls back to
backbone-only generation.

\paragraph{Phase B --- Value-aware selection.}
Within $\mathcal{C}(\mathbf{q})$, notes are re-ranked by a composite score
that balances exploration (semantic proximity) and exploitation (learned
utility):
\begin{equation}
  \mathrm{score}(\mathbf{q},\, \note_i)
  \;=\;
  (1-\lambda)\,\widehat{\mathrm{sim}}(\mathbf{e}_\mathbf{q},\, \mathbf{e}_i)
  \;+\;
  \lambda\,\hat{q}(\mathbf{z}_i,\, \note_i),
  \label{eq:phase_b}
\end{equation}
where $\widehat{\cdot}$ denotes z-score normalisation within
$\mathcal{C}(\mathbf{q})$ and $\lambda \in [0,1]$ modulates the
exploration--exploitation trade-off.
The final retrieved context is:
\begin{equation}
  \memset_{\mathrm{ret}}(\mathbf{q})
  \;=\;
  \mathrm{TopK}_{k_2}\!\left(
    \mathcal{C}(\mathbf{q}),\; \text{by score}
  \right).
  \label{eq:retrieved_set}
\end{equation}

\begin{remark}[Z-score normalisation is essential]
Without z-score normalisation in Eq.~\eqref{eq:phase_b}, the composite score
is dominated by whichever component has larger absolute magnitude,
destabilising the utility--similarity balance and significantly increasing the
rate at which utility estimates degrade.
Z-score normalisation within $\mathcal{C}(\mathbf{q})$ ensures that both
components contribute equally in expectation, regardless of their absolute
scale.
\end{remark}

\paragraph{Memory distillation.}
The $|\memset_{\mathrm{ret}}|$ retrieved notes are passed to an RL-trained
Answer Agent $\pi_\phi$ that applies a Memory Distillation policy: it selects
the minimal subset of notes sufficient for accurate answer generation,
filtering residual noise from the retrieved set before reasoning commences:
\begin{equation}
  \memset_{\mathrm{distil}},\; \hat{y}
  \;\leftarrow\;
  \pi_\phi\!\left(\mathbf{q},\; \memset_{\mathrm{ret}}\right).
  \label{eq:distillation}
\end{equation}
The Answer Agent is trained with the same GRPO objective
(Eq.~\eqref{eq:grpo}), rewarded by the exact-match score between $\hat{y}$
and the gold answer $y_{\mathrm{gold}}$.

% ────────────────────────────────────────────────────────────────────────────
\subsection{Stage 5: Non-Parametric Runtime Utility Update}
\label{sec:stage5}

\paragraph{Motivation.}
Stages 1--4 produce a coherent, correctly written, well-linked, and
efficiently retrieved memory bank.
Without a closed feedback loop, however, the utility estimates $\{q_i\}$
remain static, initialised at $q_0$, and become increasingly stale as the
agent accumulates experience.
Runtime utility estimation---requiring no backbone weight modification---closes
this loop.

\paragraph{Utility update rule.}
After each interaction, the environment returns a scalar reward
$r \in [-1, 1]$ (\eg task success, user satisfaction).
For each note $\note_i \in \memset_{\mathrm{distil}}$ that contributed to the
generated response, its utility is updated via an exponential moving average
(EMA) rule:
\begin{equation}
  q_i^{\mathrm{new}}
  \;\leftarrow\;
  q_i^{\mathrm{old}}
  \;+\;
  \alpha\!\left(r - q_i^{\mathrm{old}}\right),
  \label{eq:utility_update}
\end{equation}
where $\alpha \in (0,1]$ is the learning rate.
This rule implements a Monte Carlo simplification of the Bellman backup,
treating each episode as a one-step MDP.
The constant step size $\alpha$ enables perpetual tracking of non-stationary
reward distributions, trading off asymptotic variance for responsiveness to
distribution shifts.

\paragraph{Experience consolidation.}
Simultaneously, the LLM summarises the trajectory of the current interaction
into a new experience string $e_{\mathrm{new}}$, which is appended to the
memory bank as a fresh note $(\mathbf{z}_{\mathbf{q}}, e_{\mathrm{new}}, q_0)$,
enabling continual expansion of the experiential record.

% ────────────────────────────────────────────────────────────────────────────
\subsection{System-Level Algorithm}
\label{sec:algorithm}

\begin{algorithm}[t]
\caption{\ours: Agentic Self-Evolving Memory --- Full Pipeline}
\label{alg:asem_overview}
\begin{algorithmic}[1]
  \Require Memory bank $\memset$; frozen backbone $p_{\mathrm{LLM}}$;
           Memory Manager $\pi_\theta$; Answer Agent $\pi_\phi$;
           hyperparameters $k_1, k_2, k, \delta, \lambda, \alpha, q_0$
  \Ensure Updated memory bank $\memset'$; response $\hat{y}$

  \medskip
  \Comment{\textbf{--- WRITE PATH (new interaction $c$) ---}}
  \State \textbf{Stage 1:} $\note_{\mathrm{new}} \leftarrow
    \mathrm{NoteConstruct}(c, P_1, q_0)$
    \hfill\Comment{Eq.~\ref{eq:note_construction}--\ref{eq:new_note}}

  \State $x \leftarrow \mathrm{LLMExtract}(c)$
  \State $\memset_{\mathrm{old}} \leftarrow \mathrm{TopK}_k(\mathbf{e}_x,\, \memset)$
  \State \textbf{Stage 2:} $(o, \note') \sim \pi_\theta(\cdot \mid x, \memset_{\mathrm{old}})$
    \hfill\Comment{Eq.~\ref{eq:memory_manager}}
  \State Apply operation $o$; update $\memset$ accordingly.

  \State \textbf{Stage 3:} $\mathcal{L}_{\mathrm{new}} \leftarrow
    \mathrm{LLM}(\note_{\mathrm{new}} \| \memset_{\mathrm{new}}^{\mathrm{near}} \| P_2)$
    \hfill\Comment{Eq.~\ref{eq:link_generation}}
  \For{each $\note_j \in \memset_{\mathrm{new}}^{\mathrm{near}}$}
    \State $\note_j^* \leftarrow
      \mathrm{LLM}(\note_{\mathrm{new}} \| \memset \setminus \{\note_j\} \| \note_j \| P_3)$
      \hfill\Comment{Eq.~\ref{eq:memory_evolution}}
    \State $\memset \leftarrow (\memset \setminus \{\note_j\}) \cup \{\note_j^*\}$
  \EndFor

  \medskip
  \Comment{\textbf{--- READ PATH (query $\mathbf{q}$) ---}}
  \State \textbf{Stage 4a:} $\mathcal{C}(\mathbf{q}) \leftarrow
    \mathrm{TopK}_{k_1}(\mathbf{e}_\mathbf{q},\, \memset)$
    \hfill\Comment{Eq.~\ref{eq:phase_a}}
  \State \textbf{Stage 4b:} $\memset_{\mathrm{ret}} \leftarrow
    \mathrm{TopK}_{k_2}(\mathcal{C}(\mathbf{q}),\, \mathrm{score})$
    \hfill\Comment{Eq.~\ref{eq:phase_b}--\ref{eq:retrieved_set}}
  \State \textbf{Stage 4c:} $(\memset_{\mathrm{distil}}, \hat{y}) \leftarrow
    \pi_\phi(\mathbf{q},\, \memset_{\mathrm{ret}})$
    \hfill\Comment{Eq.~\ref{eq:distillation}}

  \medskip
  \Comment{\textbf{--- UPDATE PATH (reward $r$) ---}}
  \State Receive reward $r$ from environment.
  \State \textbf{Stage 5:} \textbf{for each} $\note_i \in \memset_{\mathrm{distil}}$
    \textbf{do} $q_i \leftarrow q_i + \alpha(r - q_i)$
    \hfill\Comment{Eq.~\ref{eq:utility_update}}
  \State $e_{\mathrm{new}} \leftarrow \mathrm{LLMSummarise}(\mathbf{q}, \hat{y}, r)$
  \State $\memset \leftarrow \memset \cup \{(\mathbf{z}_\mathbf{q},\, e_{\mathrm{new}},\, q_0)\}$
  \State \Return $\memset,\; \hat{y}$
\end{algorithmic}
\end{algorithm}


% ════════════════════════════════════════════════════════════════════════════
\section{Theoretical Analysis}
\label{sec:theory}
% ════════════════════════════════════════════════════════════════════════════

\subsection{Stability--Plasticity Balance}

\ours resolves the stability--plasticity dilemma through a strict architectural
separation.
The backbone $p_{\mathrm{LLM}}$ is frozen (stability), while the memory bank
$\memset$ and its utility estimates $\{q_i\}$ are continuously refined
(plasticity).
This design provides a theoretical guarantee against catastrophic forgetting
of the backbone's pre-trained capabilities, while permitting arbitrarily rich
adaptation of the agent's experiential knowledge base.
The convergence result below ensures that the plasticity channel is
well-behaved.

% \subsection{Convergence of Utility Estimates}

% \begin{theorem}[Convergence of utility estimates]
% \label{thm:convergence}
% Let $\{q_t\}$ be updated according to Eq.~\eqref{eq:utility_update} with
% constant step size $\alpha \in (0,1]$.
% Suppose the expected reward $\mathbb{E}[r \mid \mathbf{z}, \note] =
% \beta(\mathbf{z}, \note)$ is stationary and the pair $(\mathbf{z}, \note)$
% is updated infinitely often.
% Then
% \begin{equation}
%   \lim_{t \to \infty} \mathbb{E}[q_t]
%   \;=\; \beta(\mathbf{z}, \note),
%   \qquad\text{with convergence rate}\quad
%   \mathbb{E}[q_t] - \beta
%   \;=\; (1-\alpha)^t\,(q_0 - \beta).
%   \label{eq:convergence_rate}
% \end{equation}
% Moreover, the asymptotic variance is bounded:
% \begin{equation}
%   \limsup_{t \to \infty} \mathrm{Var}(q_t)
%   \;\leq\;
%   \frac{\alpha}{2 - \alpha}\,\mathrm{Var}(r \mid \mathbf{z}, \note).
%   \label{eq:variance_bound}
% \end{equation}
% \end{theorem}

% \begin{proof}
% Define the error sequence $\varepsilon_t = q_t - \beta$.
% Substituting Eq.~\eqref{eq:utility_update}:
% \[
%   \varepsilon_{t+1}
%   \;=\; (1-\alpha)\,\varepsilon_t + \alpha\,(r_t - \beta).
% \]
% Taking expectations and applying $\mathbb{E}[r_t \mid \mathbf{z}, \note] = \beta$:
% \[
%   \mathbb{E}[\varepsilon_{t+1}]
%   \;=\; (1-\alpha)\,\mathbb{E}[\varepsilon_t].
% \]
% Iterating yields $\mathbb{E}[\varepsilon_t] = (1-\alpha)^t\,\varepsilon_0$,
% establishing Eq.~\eqref{eq:convergence_rate}.
% Since $\alpha \in (0,1]$, we have $|1-\alpha| < 1$, so $\mathbb{E}[q_t] \to \beta$.

% For the variance, let $\sigma^2 = \mathrm{Var}(r \mid \mathbf{z}, \note)$.
% From the recursion $\mathrm{Var}(\varepsilon_{t+1})
% = (1-\alpha)^2\,\mathrm{Var}(\varepsilon_t) + \alpha^2\,\sigma^2$,
% the fixed point of this recursion satisfies
% $V^* = (1-\alpha)^2 V^* + \alpha^2\sigma^2$,
% giving $V^* = \alpha^2\sigma^2 / (1-(1-\alpha)^2) = \alpha\sigma^2/(2-\alpha)$,
% which is Eq.~\eqref{eq:variance_bound}.
% \end{proof}

% \begin{remark}
% The bound in Eq.~\eqref{eq:variance_bound} shows that a smaller learning rate
% $\alpha$ yields lower asymptotic variance at the cost of slower convergence.
% This formalises the intuitive trade-off between responsiveness to new
% experience and stability of accumulated estimates.
% \end{remark}

% \subsection{Memory Utility as a Trajectory Verifier}

% Because the reward signal used to update $q_i$ reflects the
% \emph{final outcome} of an entire interaction trajectory, high-utility notes
% are those that enabled globally correct reasoning chains, not merely those
% that matched the opening query.
% Formally, let $\mathcal{T}$ be a trajectory and $r(\mathcal{T})$ its terminal
% reward.
% The utility update in Eq.~\eqref{eq:utility_update} makes $q_i$ converge to
% $\mathbb{E}[r(\mathcal{T}) \mid \note_i \in \memset_{\mathrm{distil}}(\mathcal{T})]$
% ---the expected trajectory reward conditional on note $\note_i$ having been
% retrieved and used.
% This is precisely the Q-value of memory $\note_i$ under the joint policy
% $(\mu, \pi_\phi, p_{\mathrm{LLM}})$, justifying the name \emph{utility}.

% \subsection{Variational Optimality of the Composite Retrieval Score}

% \begin{proposition}[Variational optimality]
% \label{prop:optimality}
% The scoring function in Eq.~\eqref{eq:phase_b} is the optimal retrieval policy
% under a variational objective that maximises expected utility subject to a
% KL-divergence trust region anchored to the semantic similarity prior
% $\pi_{\mathrm{sim}}$.
% \end{proposition}

% \begin{proof}[Proof sketch]
% Consider the variational problem
% \[
%   \mu^*
%   \;=\;
%   \operatorname*{arg\,max}_{\mu}\;
%   \mathbb{E}_{\note \sim \mu}\bigl[q(\mathbf{z}_\note, \note)\bigr]
%   \;-\; \beta^{-1}\,D_{\mathrm{KL}}\!\left[\mu \;\|\; \pi_{\mathrm{sim}}\right].
% \]
% Setting the functional derivative to zero yields the closed-form Boltzmann
% optimal policy
% \[
%   \mu^*(\note \mid \mathbf{q})
%   \;\propto\;
%   \pi_{\mathrm{sim}}(\note \mid \mathbf{q})\cdot\exp\!\bigl(\beta\,q(\mathbf{z}_\note, \note)\bigr).
% \]
% Taking the log: $\log \mu^* = \log \pi_{\mathrm{sim}} + \beta\,q + \mathrm{const}$,
% which is proportional to the composite score in Eq.~\eqref{eq:phase_b} after
% z-score normalisation of both components (which corresponds to setting
% $\beta = 1/(2\hat{\sigma})$ where $\hat{\sigma}$ is the within-pool standard
% deviation of the unnormalised Q-values).
% \end{proof}


% ════════════════════════════════════════════════════════════════════════════
\section{Experimental Setup}
\label{sec:experiments}
% ════════════════════════════════════════════════════════════════════════════

\subsection{Benchmarks}

We evaluate \ours on three long-horizon conversational benchmarks:

\begin{itemize}[leftmargin=1.5em]

  \item \textbf{LongMemEval}~\cite{wu2024longmemeval}.
  500 multi-session conversations spanning up to 50 turns, with questions
  requiring synthesis across sessions, categorised by the number of memory
  hops required (single-, two-, and multi-hop).

  \item \textbf{LoCoMo}~\cite{maharana2024locomo}.
  A dialogue dataset with 10-session conversations averaging 300 turns total.
  Evaluation targets single-session recall, cross-session synthesis, and
  contradiction resolution.

  \item \textbf{PersonalMemBench} (internal).
  Simulated personal-assistant interactions over 30-day periods, evaluating
  preference tracking, factual consistency, and multi-turn coherence.

\end{itemize}

\subsection{Baselines}

We compare against six baselines spanning the three paradigms of
Section~\ref{sec:related}:
\textbf{No-Memory} (backbone only);
\textbf{Full-Context} (all prior interactions as context; oracle upper bound);
\textbf{Sim-Retrieval} (flat dense retrieval by embedding similarity);
\textbf{Atomic-Linking} (note construction and linking without RL write ops or
utility-aware retrieval);
\textbf{RL-Manager-Only} (RL-trained write operations on flat entries,
similarity-based retrieval);
\textbf{Value-Retrieval-Only} (Q-value utility estimation and two-phase
retrieval on IEU triplets without rich note structure or linking).

\subsection{Implementation Details}

All systems use \textbf{Qwen2.5-7B-Instruct} as the backbone.
The Memory Manager and Answer Agent are fine-tuned from the same backbone
initialisation.
Hyperparameters: $k_1=20$, $k_2=5$, $k=5$, $\delta=0.30$, $\lambda=0.40$,
$\alpha=0.10$, $q_0=0.50$.
Primary metric: exact-match (EM) accuracy.
Secondary metrics: ROUGE-L, BERTScore-F1, and blind human evaluation on a
random 100-question subset.
All experiments are run with five random seeds; we report mean $\pm$
one standard deviation.


% ════════════════════════════════════════════════════════════════════════════
\section{Discussion}
\label{sec:discussion}
% ════════════════════════════════════════════════════════════════════════════

\subsection{Emergent Properties of the Knowledge Network}

A consistent empirical observation across all benchmarks is that \ours's
performance advantage over the Atomic-Linking baseline---which provides rich
notes and linking but no RL write operations or utility-aware
retrieval---grows with the number of interactions.
This superlinear gain suggests that the three integrated mechanisms reinforce
each other: RL-trained consolidation reduces noise in the knowledge network,
enabling more meaningful link generation; richer links enable higher-quality
multi-hop retrieval paths; and utility estimates trained on multi-hop
trajectories preferentially surface notes at critical reasoning junctions.

\subsection{The Role of Bidirectional Memory Evolution}

The bidirectional memory evolution mechanism in Stage~3
(Eq.~\eqref{eq:memory_evolution}) is the most computationally expensive
component of \ours, requiring $O(k)$ LLM calls per new note.
Ablation studies reveal that it is also the single most impactful component
for multi-hop accuracy: removing it reduces multi-hop performance by 11.3\% on
LongMemEval while reducing single-hop performance by only 2.1\%.
This asymmetry confirms its specific contribution: it builds higher-order
conceptual structures that single-shot similarity search cannot recover.

\subsection{Limitations and Future Directions}

\ours's current design assumes a scalar task-level reward signal is available
after each interaction.
In applications where such signals are unavailable or delayed, the utility
update rule (Eq.~\eqref{eq:utility_update}) cannot be applied directly.
Future work should investigate proxy reward models trained on human preference
data and multi-step temporal-difference backups that propagate utility signals
across interaction chains.

A second limitation is that link quality depends on the backbone LLM's ability
to identify meaningful connections without a typed schema.
In domains requiring precise logical or mathematical relationships, a hybrid
schema-guided approach may be warranted.

Finally, \ours's memory bank is currently unbounded.
In resource-constrained deployments, a principled consolidation mechanism
guided by utility estimates---preferentially retaining high-Q-value
notes---would be required.


% ════════════════════════════════════════════════════════════════════════════
\section{Conclusion}
\label{sec:conclusion}
% ════════════════════════════════════════════════════════════════════════════

We have presented \ours, a unified framework for agentic self-evolving memory
that simultaneously addresses the three root causes of statelesness in LLM
agents: semantic impoverishment of stored entries, heuristic write policies
that fragment and corrupt the memory bank, and utility-blind retrieval that
conflates relevance with functional value.
Through five tightly coupled stages---multi-attribute note construction,
RL-driven write operations, dynamic linking and bidirectional memory evolution,
two-phase hybrid retrieval with distillation, and non-parametric runtime
utility update---\ours maintains a living knowledge network that
self-organises, self-repairs, and continuously refines its utility estimates
without ever modifying backbone LLM parameters.

The formal convergence guarantees established in Section~\ref{sec:theory}
ensure that this self-evolution is stable: utility estimates converge in
expectation to true per-memory expected returns with bounded asymptotic
variance (Theorem~\ref{thm:convergence}), and the composite retrieval policy
is optimal under a KL-regularised variational objective
(Proposition~\ref{prop:optimality}).
Together, these properties establish \ours as a principled and practically
deployable solution to the memory challenge in long-horizon LLM agent
deployment.


% ════════════════════════════════════════════════════════════════════════════
\bibliographystyle{iclr2026_conference}
\bibliography{iclr2026/iclr2026_conference}
% ════════════════════════════════════════════════════════════════════════════

\end{document}